\subsection{Estructura de directorios}

El Código~\ref{lst:estructura} muestra la organización de carpetas del
código fuente y su correspondencia con las capas de la arquitectura.

\begin{lstlisting}[language={}, numbers=none, caption={Estructura de carpetas del código fuente}, label={lst:estructura}]
coach-app/
  src/
    app/                  VISTA: rutas de Expo Router
      _layout.tsx         rutas protegidas segun sesion y perfil
      index.tsx           bienvenida
      (auth)/             login, register
      (onboarding)/       goal, physical-data, routine
      (app)/              pestanas de la app autenticada
        index.tsx         inicio
        routines/         index, generate, custom, exercise (ficha)
        machines/         index, [id] (ficha), scanner
        nutrition.tsx     diario de alimentacion y calculadora de IMC
        progress/         peso e IMC, entrenos y reporte en PDF
        profile.tsx       perfil (se abre desde el avatar)
        train.tsx         entrenamiento guiado (pantalla completa)
    components/           VISTA: componentes reutilizables por modulo
    controllers/          CONTROLADOR: Auth, Machines, Nutrition, Progress,
                          Reports, Routines, WgerExercises, Workouts
    utils/                CONTROLADOR: reglas puras (bmiCalculator,
                          nutritionGoals, routineGenerator, workoutStats,
                          progressReport, foodTranslator, ...)
    hooks/                logica de pantalla y suscripciones en tiempo real
    store/                estado global (Zustand): auth, onboarding, toast
    models/
      entities/           MODELO: tipos del dominio
      repositories/       MODELO: acceso a datos por coleccion
    schemas/              MODELO: validacion con Zod
    services/
      firebase/           SERVICIOS: Firestore, Storage, Functions, App Check
      wger/  usda/        SERVICIOS: wger (HTTP) y USDA (via funcion)
      exercisedb/         SERVICIOS: demostraciones alternativas
      ai/                 SERVICIOS: TensorFlow Lite local y deteccion en la
                          nube (routinePlanner, machineRecognition)
      ar/                 SERVICIOS: interfaz tipada de la vista CoachAR
      reports/            SERVICIOS: reporte de progreso en PDF
    config/  constants/   entorno y Google Sign-In; paleta de colores
    types/  media/        tipos declarados; GIF de demostracion locales
    **/__tests__/         pruebas Jest junto al codigo de cada carpeta
  modules/coach-ar/       modulo nativo de ARCore (Kotlin): vista CoachAR
  assets/models/          modelos .tflite (rutinas y reconocimiento)
  android/                proyecto nativo de Android
  ml/                     entrenamiento de los modelos (Python) y metricas
    detection/            detector de maquinas YOLO: datos y evaluacion
    recognition-yolo/     clasificador de maquinas (YOLO11n-cls)
  functions/              Cloud Functions (Node.js 22): buscarAlimentos,
                          guardarRutina, reconocerMaquina, cuotas
  tests/security/         pruebas de las reglas con el emulador de Firebase
  docs/                   ARCORE.md, auditorias y despliegue del backend
  firestore.rules, storage.rules, firebase.json   reglas y configuracion
  .github/workflows/      integracion continua (backend-security.yml)
\end{lstlisting}
